\gls{ELF} is the standard executable format under Linux operating systems.
The work done on this thesis heavily involve the \gls{ELF} structure, with particular focus on some sections 
used for relocation, alongside the process of execute an \gls{ELF}.
\subsection{ELF Structure}
    In this subsection the process of \gls{ELF} execution is discussed with a focus on the parts of the execution flow 
    that will influence this work.\\
    \textbf{ELF Header} is a structure locate at the beginning of the \gls{ELF} and contains some important informations 
    about the type of \gls{ELF}, the type of target operating system \dots.
    Alongside these informations there is the location of two important data structure inside the \gls{ELF} :
    \begin{itemize}
        \item The \textbf{Program Header Table} that defines the \textit{Execution View}
        \item The \textbf{Section Header Table} that defines the \textit{Linking View}
    \end{itemize}
    The \textbf{Program Header Table} contains informations to create a runtime image of the program in memory with 
    the  \textbf{Segments} \textit{(i.e. which parts of the fils must be loaded into memory (), which shared 
    objects are needed \dots)}.
    Each segment contains the access flag for that part of memory (i.e. Read/Write/Execute).
    The \textbf{Section Header table} contains the informations about the \textbf{Sections} that compose the binary, 
    togheter with the mapping with the segments (the section \textit{.text} in which segment is mapped?).
    Among all the sections, some of them contains specific structure used to perform specific tasks:
    \begin{itemize}
        \item \textit{\textbf{.dynamic} : The dynamic section contains informations used by the Dynamic Linker 
            (e.g. the dependencies, the interpreter location \dots)}.
        \item \textit{\textbf{.dynsym} : List of symbols defined and used by the binary. 
            If the binary is not stripped, there is also the \textbf{.symtab} section, that contains a more complete list of symbols}.
        \item \textit{\textbf{.hash or .gnu.hash} : Structures used for efficient symbol lookup during symbol resolution}\footnote{
            The \textbf{.gnu.hash} section is the more efficient version of \textbf{.hash}
        }.
        \item \textit{\textbf{.rela.dyn and .rela.plt} : Both are relocation table. the first is the standard one, 
        the second is used in case of \textbf{Lazy Binding}}.
        \item \textit{\textbf{.shstrtab, .dynstr and .strtab} : Holds lists of strings referenced by other parts of the \gls{ELF}.
            The first holds section names, the second holds the names of the symbols contained into \textbf{.dynsym} section. 
            the last one holds the names of the entire \textbf{.symtab} section}.
        \item \textit{\textbf{.got and .got.plt} : Contains the \gls{GOT} filled by the dynamic linker with a
            ddresses of external functions and variables. The first is the standard one, the second is used by the \gls{PLT}}.
        \item \textit{\textbf{.plt} : Contains the \gls{PLT} that is an executable section with trampolines. In particular 
            contains stubs for function that might be located in another \gls{ELF} or \gls{SO}. 
            The stub reference the corresponding \gls{GOT} entry with the resolved address of the function.
            This section is used in case of \textbf{Lazy Binding}}.
        \item \textit{\textbf{.init and .init\_array} : Sections used for initialization. \textbf{.init} contains initialization code.
            \textbf{.init\_array} contains addresses of initialization functions}.
        \item \textit{\textbf{.fini and .fini\_array} : Sections used for termination handling. 
            \textbf{.init} contains termination code.
            \textbf{.init\_array} contains addresses of termination functions}.
    \end{itemize}
\subsection{ELF Linking}
    Before actually dig deep into \gls{ELF} loading, the linking process must be revised.
    The \textbf{Linking phase} consist in join togheter different pieces of code belonging to different files.
    For example if a program uses \texttt{printf} function, it need to take the code of the printf from the libc library.
    This process of joinig togheter the executable that uses the printf with the printf code that is in another file 
    is called \textbf{Linkinig}. The linking could be of two different types : 
    \begin{itemize}
        \item \textbf{Static linking} : The needed code (e.g. the instructions of printf) is copied inside the binary 
        at \textit{Link Time}. This cause the dimension of the binary to be bigger (because the code 
        of external functions actually lies inside it) and conseguently a waste of RAM space (if 10 binaries both uses 
        printf function, the code of the printf is copied inside all the 10 binaries).
        \item \textbf{Dynamic linking} : The needed code is not copied inside the binary, but it contains informations (i.e. 
        the address of the external function or variable) that are used to locate the needed resource.
        The resolution happens at runtime when the Dynamic Linker finds and load in the process virtual address space the 
        library and resolve the reference to the symbols. This allows different processes to share the same \gls{SO}.
    \end{itemize}
    Since nowaday the most used approach is \textbf{Dynamic Linking} the following discussion will focus on it.\\
    As stated before the linking process can be used for both external variables and external functions.
    In both cases the process is the same since both functions and variables could be exported as symbols.\\
    When dealing with functions there is the possibility of use \textbf{Lazy Binding}.
    This technique consist in resolve the symbol only during the first invocation of the function and not at load time 
    as happens with the variables.
    In order to achieve \textbf{Lazy Binding} a stub section is used (i.e. \textbf{.plt}).
    When a program uses a library function, a \gls{PLT} entry is added into the .plt section.
    \begin{figure}[ht]
    \centering
            \begin{subfigure}[h]{0.7\textwidth}
            \centering
            \includegraphics[width=\linewidth]{LazyBindingExample.png}
            \caption{An example of puts function resolution using Dynamic Linking and Lazy Binding before the actual resolution}
            \label{fig:lazy_resolution_plt_1}
        \end{subfigure}
        \vspace{0.5cm}
        \begin{subfigure}[h]{0.7\textwidth}
            \centering
            \includegraphics[width=\linewidth]{LazyBindingExampleAfterResolution.png}
            \caption{Same binary of \cref{fig:lazy_resolution_plt_1} but with symbol resolved}
            \label{fig:lazy_resolution_plt_2}
        \end{subfigure}
        \caption{An example Lazy Binding involving .plt and .got}
        \label{fig:lazy_resolution}
    \end{figure}
    As can be seen in \cref{fig:lazy_resolution} the relocation is resolved lazily upon the first invocation of the function. 
    When control is transferred to the corresponding \gls{PLT} entry, the entry performs an indirect jump through 
    its associated \gls{GOT} entry. Before the symbol is resolved, this \gls{GOT} entry contains an address that 
    redirects execution back to the \gls{PLT}, eventually reaching the \gls{PLT} resolver stub. 
    The resolver then invokes the dynamic linker, which performs the symbol lookup, resolves the relocation, 
    and updates the \gls{GOT} entry with the resolved address of the function as can be seen in \cref{fig:lazy_resolution_plt_2}. 
    Subsequent calls can therefore jump directly to the resolved address through the \gls{GOT}.

\subsection{ELF Loading}
    The next useful step in order to understand the choosen approach consist in go into details of ELF Loading phase.
    There are 3 main actors that allow a binary to be executed properly:
    \begin{itemize}
        \item \textbf{Kernel} : is the one responsible of the actual \texttt{execve} invocation and 
            allocation of different segments of memory. Moreover is the one that execute 
            the dynamic linker used to resolve dependencies prior the execution of the binary
        \item \textbf{Dynamic Linker} : Loads all the needed libraries into the process memory, 
            resolve the external dependencies, prepare the .got and .plt and execute the initializers of the libraries.
        \item \textbf{Target Binary} : Invoked by the Dynamic Linker, it prepare the environment for the main function and, 
            in the end, it runs it.
    \end{itemize}

    \begin{figure}[ht]
        \centering
        \includegraphics[width=300pt]{dynamicallyLinkedProgramLoad.png}
        \caption{Dynamically linked ELF loading procedure}
        \label{fig:dynamicallyLinkedProgramLoad}
    \end{figure}
    All these procedures could be seen in details in \cref{fig:dynamicallyLinkedProgramLoad}.
    The dynamic linker part is the one on which I will focus in this dissertation.
    The dynamic linker:
    \begin{itemize}
        \item \textit{Loads all the DT\_NEEDED entries} : these are the dependencies (i.e. \gls{SO} needed by the binary).
            For example if the binary uses an external function contained in \textit{libc.so}, the library from which the 
            function is exported must be added as dependency
        \item \textit{Resolve the relocations} : a focus on relocations will be presented in the next subsection
        \item \textit{Arranges \gls{GOT} and \gls{PLT}}
        \item \textit{Execute all the initializer functions (constructor) for each dependency} : For each \textit{DT\_NEEDED} 
            entry found by the dynamic linker, it execute all the functions contained in \textit{.init\_array} section of 
            that entry. These are functions declared with \texttt{\_\_attribute\_\_((constructor))}.
        \item \textit{Find and execute the binary entry point}
    \end{itemize}

\subsection{ELF Symbols and Relocations}
    The dynamic linker will use the symbols to link different ELF togheter.
    \begin{itemize}
        \item A component (an \gls{ELF} or a \gls{SO}) could export a symbol that represent a variable or a function inside it.
        \item Another component could reference the exported symbol and use it (if the first component is a dependency of the second).
        \item The linker is the intermediary between the two components that is responsible of match the symbols names 
            (the exported and the referenced) and store the address of exported symbol into the second component address space 
            in a way that this component could access it.
    \end{itemize}
    All the symbols that belong to an \gls{ELF} (both imported and exported) are stored into \textit{.dynsym} section.
    \begin{figure}[ht]
        \centering
        \includegraphics[width=\linewidth]{dynsym_example.png}
        \caption{Example of a \textit{.dynsym} table for a shared object}
        \label{fig:dynsym_example}
    \end{figure}
    \begin{figure}[ht]
        \centering
        \includegraphics[width=\linewidth]{reladynexample.png}
        \caption{Example of a \textit{.rela.dyn} relocation table extracted from a binary}
        \label{fig:reladyn_example}
    \end{figure}
    In \cref{fig:dynsym_example} can be seen a list of symbols extracted from a \gls{SO}.
    The first 4 symbols are external symbols that must be resolved. This can be seen because the Value column is 0 and Ndx is UND (undefined).
    The last 7 instead are symbols declared inside the shared object.
    When the visibility is setted to \textit{DEFAULT} this means that the symbol can be visible outside the shared object.

    The process of resolving the symbol references is called \textbf{Relocation}. The relocation algorithm is shown in 
    \cref{alg:relocation_resolution}.
    Relocate a symbol means resolve the symbol (i.e. find the address of that symbol in the original object), and write 
    that address inside the binary that has the relocation.
    Specifically the symbol address is written inside the \gls{GOT}.
    In \cref{fig:reladyn_example} the address where the symbol address must be written can be seen in the ``Offset'' column.
    \begin{algorithm}
        \caption{Relocation resolution at ELF load time}
        \label{alg:relocation_resolution}
        \begin{algorithmic}[1]
            \REQUIRE $L = [obj_1,\dots,obj_n]$ \COMMENT{The object that has a dependency (DT\_NEEDED) must already been loaded in memory}
            \ENSURE for each $obj \in L$, the respective GOT entry contains the resolved runutime address of the external symbol.
            
            \FOR{$obj \in L$}
                \STATE $Reloc\_table \gets obj.\texttt{GET\_SECTION(".rela.dyn")}$
                \FOR{$reloc \in Reloc\_table$}
                    \IF{$reloc.\text{type} = \texttt{R\_*\_RELATIVE}$}
                        \STATE $addr \gets obj.\texttt{loadBias} + r.\text{addend}$ 
                        \COMMENT{\textcolor{green!60!black}{The symbol does not need to be resolved. The address must be adjusted in order to enable ASLR}}
                    \ELSE
                        \IF{$reloc.\text{type} \in \{\texttt{R\_*\_GLOB\_DAT} , \texttt{R\_*\_JUMP\_SLOT}\}$} 
                            \STATE $sIdx \gets reloc.\text{symbol\_index}$
                            \STATE $symbolName \gets obj.\texttt{dynstr}[\, obj.\texttt{dynsym}[sIdx].\text{name} \,]$ 
                            \COMMENT{\textcolor{green!60!black}{The name of the symbol is read from \textit{.dynsym}}}
                            \STATE $target \gets \textsc{ResolveSymbol}(symbolName, L, obj)$
                            \IF{$target = \text{NULL}$}
                                \STATE \textbf{error}: undefined symbol $symbolName$ 
                            \ELSE
                                \STATE $(defObj, defIdx) \gets target$ \COMMENT{\small \textcolor{green!60!black}{\textit{defObj} is the object that export the symbol, \textit{defIdx} is the index of the .dynsym table of \textit{defObj}}}
                                \STATE $addr \gets defObj.\texttt{loadBias} + defObj.\texttt{dynsym}[defIdx].\text{value}$
                                \COMMENT{\small \textcolor{green!60!black}{The address of the resolved symbol is computed as base\_address (of the object that export the symbol) + the offset in which the symbol is stored}}
                            \ENDIF
                        \ENDIF
                    \ENDIF
                    \STATE write $addr$ in $obj.\texttt{got}$ at offset $reloc.\text{offset}$ \COMMENT{\small \textcolor{green!60!black}{The destination address in which write the resolved symbol address is specified in the relocation entry}}
                \ENDFOR
            \ENDFOR
        \end{algorithmic}
        
    \end{algorithm}
    \begin{algorithm}
        \caption{\textsc{ResolveSymbol} function used by \cref{alg:relocation_resolution}}
        \label{alg:symbol_resolution}
        \begin{algorithmic}[1]
            \REQUIRE name $name$, object list $L$, object that request the resolution $requester$
            \ENSURE pair $(obj, idx)$ of the first found symbol, or NULL if the symbol has not been defined by any object.
            \FOR{$obj \in L$ \COMMENT{\small \textcolor{green!60!black}{The order is: first the executable, then all the DT\_NEEDED dependencies in order of definition}}}
                \STATE $idx \gets \textsc{GnuHashLookup}(name,\ obj.\texttt{gnuHash},\ obj.\texttt{dynsym},\ obj.\texttt{dynstr})$
                \COMMENT{\small \textcolor{green!60!black}{As seen before, the hashes are used in order to speedup the symbol lookup}}
                \IF{$idx \neq \text{NOT\_FOUND}$ \textbf{\&\&} $obj.\texttt{dynsym}[idx].\text{st\_shndx} \neq \texttt{SHN\_UNDEF}$}
                    \STATE \textbf{return} $(obj, idx)$ 
                \ENDIF
            \ENDFOR
            \STATE \textbf{return} NULL
        \end{algorithmic}
    \end{algorithm}
